\clearpage
\hypertarget{native-483-parthenope}{}
\begin{center}
\rule{0.55\linewidth}{0.4pt}\\[2pt]
Astronomische Nachrichten. Band XXX. Nr. 719. S. 374. 20. Juni 1850.\\[2pt]
\rule{0.55\linewidth}{0.4pt}
\end{center}
\begin{center}Beobachtungen und Elemente der Parthenope.\end{center}

Geheimrat \textsc{Gauss} hat Parthenope am Meridiankreis der Göttinger Sternwarte beobachtet,

\begin{center}
\setlength{\tabcolsep}{4pt}
\begin{tabular}{l|r|r|r}
 & Mittlere Zeit & Rektasz. & Südl. Dekl. \\\hline
Mai 31 & $10^{\rm u}\ 27^{\rm m}\ 38^{\rm s}$ & $225^\circ\ 52'\ 38''4$ & $9^\circ\ 48'\ 3''7$ \\
\end{tabular}
\end{center}

Die Luft war in den Schichten näher am Horizont etwas trüb. Bei schwächster Beleuchtung der Fäden waren die Durchgänge nur schwer zu erkennen.

\begin{center}
\setlength{\tabcolsep}{4pt}
\begin{tabular}{l|r|r|r}
Juni 2 & $10^{\rm u}\ 18^{\rm m}\ 1^{\rm s}6$ & $225^\circ\ 30'\ 3''6$ & $9^\circ\ 45'\ 59''8$ \\
\end{tabular}
\end{center}

